\begin{proof}Keeping the established notation, it suffices to show that, if~$Z(p)$ has rank at least~$2$, then there exists a $P\in\GL(3,\bbR)$ such that, when $Q$ is defined by~\eqref{eq: QintermsofPandZ}, the tableau~$A_Q$ is involutive.

Now, by the hypothesis that there is no nonzero vector $\mbv\in T_pM$
such that $\mbv\lhk\Omega=0$, the rank of $Z(p)$ is either $2$ or $3$.
When the rank of $Z(p)$ is~$3$,
as $P$ varies over~$\GL(3,\bbR)$,
the matrix~$Q$ varies over an open subset of~$\GL(3,\bbR)$,
and it is clear that, for the generic choice of~$P$,
the corresponding $Q_0$ will not be $\SO(3)$-equivalent to anything
in the two `degenerate' cones defined by~\eqref{eq: Q0normalizeddegenerate}
with either $r=y=0$ or $r^2=x^2+y^2$.

When the rank of $Z(p)$ is $2$, we can assume, after an $\SO(3)$ rotation,
that the bottom row of~$Z(p)$ vanishes and that the first two rows of~$Z(p)$
are linearly independent. It then follows that $P/(\det P)\, {}^t\!Z(p)$
has its last column equal to zero, but that, as $P$ varies, the first
two columns of $P/(\det P)\, {}^t\!Z(p)$ range over all linearly independent
pairs of column vectors. Now explicitly computing the polynomial $c_Q(\mbz^*)$
for the corresponding matrix~$Q$ shows that $c_Q(\mbz^*)$ does not vanish
identically on the set of such matrices, hence it is possible to choose $P$
so that $c_Q(\mbz^*)$ does not vanish identically, and the corresponding $A_Q$
is then involutive, implying that the corresponding admissible integral
element~$E$ is Cartan-ordinary.

In either case, there exist Cartan-ordinary admissible integral elements
of~$\cI$ based at~$p$, so the Cartan--K\"ahler theorem applies, showing
that there exist admissible integral manifolds of~$\cI$ passing through
any point~$(p,\mbx,\mba)\in X^9$, and hence,
by Proposition~\ref{prop: formulation_as_EDS},
the original problem is solvable in an open neighborhood of~$p$.
Moreover, since the last nonzero Cartan character
of a generic integral flag is~$s_2 = 3$,
the space of solutions~$\omega$
depends locally on $3$ functions of $2$ variables,
in the sense of Cartan.
\end{proof}
